\documentclass[aspectratio=169]{beamer}
\usetheme{Madrid}
\usecolortheme{default}
\usepackage[utf8]{inputenc}
\usepackage{amsmath}
\usepackage{amsfonts}
\usepackage{amssymb}
\usepackage{enumerate}

\setbeamertemplate{navigation symbols}{}
\setbeamerfont{block body}{size=\small}
\setbeamerfont{alertblock body}{size=\small}

\title{CRISC Mixed Practice Exam -- 100 Real-Life Questions}
\subtitle{Weighted to the Four CRISC Domains -- Practical Scenarios}
\author{CRISC Mastery Series}
\date{\today}

% ------------------------------------------------------------
% \hq{number}{topic}{question}{options}{letter}{explanation}{takeaway}
% Creates one question frame and one explanation frame.
% ------------------------------------------------------------
\newcommand{\hq}[7]{%
	\begin{frame}{Q#1: #2}
		\begin{block}{Question}
			#3
		\end{block}
		\begin{itemize}
			#4
		\end{itemize}
	\end{frame}
	\begin{frame}{Q#1: Explanation}
		\begin{block}{Correct Answer: #5}
			\begin{itemize}
				#6
			\end{itemize}
		\end{block}
		\begin{alertblock}{Key Takeaway}
			#7
		\end{alertblock}
	\end{frame}
}

\begin{document}
	
	\begin{frame}
		\titlepage
	\end{frame}
	
	\begin{frame}{Instructions}
		\begin{alertblock}{How to Use This Deck}
			Think like a risk practitioner in a real organization.\\
			Focus on business context, ownership, and practical decision-making.\\
			The answer and a key takeaway are on the slide following each question.
		\end{alertblock}
	\end{frame}
	
	\begin{frame}{Outline}
		\tableofcontents
	\end{frame}
	
	% ============================================================
	\section{Governance (26 percent of the exam)}
	% ============================================================
	
	\hq{1}{Risk Committee Charter}
	{A newly formed risk committee has no written charter. What is the \textbf{MOST} important item to define first?}
	{\item[\textbf{(A)}] The brand of meeting software
		\item[\textbf{(B)}] Its purpose, authority, membership, and reporting line to the board
		\item[\textbf{(C)}] The list of approved vendors
		\item[\textbf{(D)}] The committee's lunch schedule}
	{B}
	{\item A charter establishes legitimacy and scope.
		\item Authority determines what decisions it can make.
		\item A clear reporting line supports oversight.}
	{Define authority before activity.}
	
	\hq{2}{Enterprise Risk Integration}
	{IT risk is managed separately from enterprise risk, producing conflicting ratings. What is the \textbf{BEST} improvement?}
	{\item[\textbf{(A)}] Keep them separate to preserve independence
		\item[\textbf{(B)}] Let IT use its own scale permanently
		\item[\textbf{(C)}] Remove IT risk from reports
		\item[\textbf{(D)}] Integrate IT risk into the enterprise risk framework using common criteria and reporting}
	{D}
	{\item Common criteria allow comparison across risk types.
		\item IT risk affects business objectives.
		\item Integrated views support better decisions.}
	{IT risk is business risk.}
	
	\hq{3}{Risk Taxonomy}
	{Departments describe the same risk using different terms, causing confusion in reports. What should be implemented?}
	{\item[\textbf{(A)}] A common risk taxonomy and glossary approved by governance
		\item[\textbf{(B)}] More frequent meetings with no changes
		\item[\textbf{(C)}] A larger spreadsheet
		\item[\textbf{(D)}] A rule that only IT may name risks}
	{A}
	{\item Shared language reduces misinterpretation.
		\item Governance approval gives it authority.
		\item Consistency improves aggregation.}
	{Speak one risk language.}
	
	\hq{4}{Vendor Gifts}
	{A vendor under evaluation offers the risk practitioner expensive concert tickets. What is the \textbf{BEST} response?}
	{\item[\textbf{(A)}] Accept because it is only a gift
		\item[\textbf{(B)}] Accept and rate the vendor higher
		\item[\textbf{(C)}] Decline and disclose the offer in line with the organization's code of conduct
		\item[\textbf{(D)}] Accept and tell nobody}
	{C}
	{\item Gifts can influence objectivity.
		\item Policies usually require disclosure.
		\item Independence protects the assessment.}
	{Protect objectivity; decline and disclose.}
	
	\hq{5}{Incentives and Risk Behavior}
	{Sales bonuses reward revenue only, and staff bypass security approvals to close deals. What is the \textbf{ROOT} governance issue?}
	{\item[\textbf{(A)}] Salespeople are inherently careless
		\item[\textbf{(B)}] Incentives are not aligned with the organization's risk appetite and policies
		\item[\textbf{(C)}] Security approvals are too short
		\item[\textbf{(D)}] Software is too old}
	{B}
	{\item People respond to what is rewarded.
		\item Governance should align incentives with risk behavior.
		\item Fixing symptoms without incentives will not hold.}
	{Culture follows incentives.}
	
	\hq{6}{IT Steering Committee}
	{Which role is \textbf{MOST} appropriate for an IT steering committee?}
	{\item[\textbf{(A)}] Configure servers
		\item[\textbf{(B)}] Write all code
		\item[\textbf{(C)}] Prioritize IT investments and monitor alignment with business strategy
		\item[\textbf{(D)}] Handle help desk tickets}
	{C}
	{\item Steering committees provide direction and oversight.
		\item Business and IT leaders both participate.
		\item Operational tasks stay with management.}
	{Steering sets direction, not settings.}
	
	\hq{7}{Delegation of Authority}
	{A project manager approves a risk acceptance far above their delegated limit. What is the \textbf{BEST} action?}
	{\item[\textbf{(A)}] Accept it because it is already signed
		\item[\textbf{(B)}] Treat the approval as invalid and route it to the appropriate authority
		\item[\textbf{(C)}] Hide the document
		\item[\textbf{(D)}] Increase the limit retroactively without review}
	{B}
	{\item Authority limits control who may accept risk.
		\item Invalid approvals leave risk unmanaged.
		\item Retroactive changes undermine governance.}
	{Approval is valid only at the right level.}
	
	\hq{8}{Policy Conflicts}
	{Two policies give conflicting retention requirements for the same data. What should the risk practitioner do?}
	{\item[\textbf{(A)}] Follow whichever is easier
		\item[\textbf{(B)}] Ignore both policies
		\item[\textbf{(C)}] Raise the conflict with policy owners and legal so governance can resolve and harmonize it
		\item[\textbf{(D)}] Delete the data}
	{C}
	{\item Conflicts create compliance risk.
		\item Policy owners must reconcile them.
		\item Legal input confirms obligations.}
	{Resolve policy conflicts at the source.}
	
	\hq{9}{Regulatory Change Management}
	{New regulations are published several times a year. Which governance practice is \textbf{MOST} effective?}
	{\item[\textbf{(A)}] Wait for audits to mention changes
		\item[\textbf{(B)}] Assign changes to whoever is available
		\item[\textbf{(C)}] A defined process to monitor, interpret, and assign actions for regulatory changes
		\item[\textbf{(D)}] Review regulations every five years}
	{C}
	{\item Monitoring prevents surprise non-compliance.
		\item Assigned owners ensure action.
		\item Process makes response repeatable.}
	{Make regulatory change a process, not a surprise.}
	
	\hq{10}{Data Governance Council}
	{What is the \textbf{PRIMARY} benefit of a data governance council?}
	{\item[\textbf{(A)}] Faster hardware purchases
		\item[\textbf{(B)}] Consistent ownership, classification, quality, and usage decisions across the enterprise
		\item[\textbf{(C)}] Lower electricity costs
		\item[\textbf{(D)}] Fewer staff meetings}
	{B}
	{\item Data decisions cross business units.
		\item A council clarifies owners and rules.
		\item Quality and privacy risks fall.}
	{Shared data needs shared rules.}
	
	\hq{11}{Risk-Based Audit Plan}
	{How should internal audit \textbf{BEST} determine its annual audit plan?}
	{\item[\textbf{(A)}] Audit the same areas each year
		\item[\textbf{(B)}] Audit by alphabetical department
		\item[\textbf{(C)}] Let each department choose
		\item[\textbf{(D)}] Using a risk-based approach that reflects the organization's key risks and objectives}
	{D}
	{\item Risk-based planning focuses on what matters.
		\item Plans should respond to changing risk.
		\item Coverage should be explainable to the board.}
	{Audit where the risk is.}
	
	\hq{12}{Merger and Appetite}
	{After acquiring a company with a very different risk culture, what should the board do regarding risk appetite?}
	{\item[\textbf{(A)}] Leave it unchanged forever
		\item[\textbf{(B)}] Adopt the smaller company's appetite without discussion
		\item[\textbf{(C)}] Review and, where needed, update the enterprise appetite and communicate it to the combined organization
		\item[\textbf{(D)}] Eliminate appetite statements}
	{C}
	{\item Business change can alter risk exposure.
		\item Appetite should reflect current strategy.
		\item Communication drives consistent behavior.}
	{Re-examine appetite after major change.}
	
	\hq{13}{Project Risk Roles}
	{A large project has no defined risk roles. What is the \textbf{BEST} action?}
	{\item[\textbf{(A)}] Wait until problems occur
		\item[\textbf{(B)}] Make the vendor responsible
		\item[\textbf{(C)}] Define risk roles and responsibilities, including who owns and reviews project risks
		\item[\textbf{(D)}] Assign all risks to the project sponsor's assistant}
	{C}
	{\item Clear roles ensure risks are managed.
		\item Owners need authority.
		\item Early definition prevents gaps.}
	{Define who does what before the project starts.}
	
	\hq{14}{Shadow IT Governance}
	{Business units increasingly buy cloud tools without IT involvement. What is the \textbf{BEST} governance response?}
	{\item[\textbf{(A)}] Provide an efficient, risk-based intake process and clear policy so teams can adopt tools safely
		\item[\textbf{(B)}] Ban all cloud tools
		\item[\textbf{(C)}] Ignore it
		\item[\textbf{(D)}] Punish all teams}
	{A}
	{\item Friction drives shadow IT.
		\item A simple process improves visibility.
		\item Risk-based review keeps oversight proportionate.}
	{Make the secure path the easy path.}
	
	\hq{15}{Risk Awareness Training}
	{Which measure \textbf{BEST} indicates that risk awareness training is working?}
	{\item[\textbf{(A)}] The number of slides in the course
		\item[\textbf{(B)}] Improved behavior, such as lower phishing click rates and more timely issue reporting
		\item[\textbf{(C)}] Hours of video watched
		\item[\textbf{(D)}] The number of emails sent}
	{B}
	{\item Outcomes show real change.
		\item Completion rates alone are weak.
		\item Behavior metrics guide improvement.}
	{Measure behavior, not attendance.}
	
	\hq{16}{Board Reporting Cadence}
	{The board receives risk reports only once a year. What is the \textbf{MOST} significant concern?}
	{\item[\textbf{(A)}] Reports are too short
		\item[\textbf{(B)}] The board has too little coffee
		\item[\textbf{(C)}] Emerging risks and breaches of appetite may not be reported in time for oversight
		\item[\textbf{(D)}] Reports are in color}
	{C}
	{\item Oversight needs timely information.
		\item Material changes require escalation outside the cycle.
		\item Frequency should reflect volatility.}
	{Report at the speed of risk.}
	
	\hq{17}{Outsourcing Governance}
	{The organization outsources a critical process. Which governance element is \textbf{MOST} important?}
	{\item[\textbf{(A)}] Only the lowest price
		\item[\textbf{(B)}] Trusting verbal promises
		\item[\textbf{(C)}] Removing internal oversight
		\item[\textbf{(D)}] Clear service levels, defined oversight responsibilities, and exit and continuity provisions}
	{D}
	{\item Oversight must be defined internally.
		\item Exit provisions reduce lock-in.
		\item Service levels make performance measurable.}
	{Outsource the work, keep the oversight.}
	
	\hq{18}{Ethical AI Use}
	{The company deploys AI for hiring screening. Which governance step is \textbf{MOST} important?}
	{\item[\textbf{(A)}] Let the algorithm decide alone
		\item[\textbf{(B)}] Publish no information
		\item[\textbf{(C)}] Establish accountability, bias testing, transparency, and human review of decisions
		\item[\textbf{(D)}] Use only the cheapest model}
	{C}
	{\item AI can embed bias.
		\item Humans remain accountable.
		\item Transparency supports trust.}
	{Automated decisions need human accountability.}
	
	\hq{19}{Policy Exception Trends}
	{Exception requests for one policy have tripled in a year. What should the risk practitioner conclude and do?}
	{\item[\textbf{(A)}] The policy may be impractical or misaligned, so review it with its owner and the business
		\item[\textbf{(B)}] Approve all exceptions automatically
		\item[\textbf{(C)}] Stop accepting exception requests
		\item[\textbf{(D)}] Conclude users are malicious}
	{A}
	{\item Patterns of exceptions signal design problems.
		\item Policy owners should review applicability.
		\item Governance should adjust or enforce deliberately.}
	{Many exceptions mean the policy needs attention.}
	
	\hq{20}{Business Ownership of IT Risk}
	{IT is asked to accept risks arising from a business-owned application. What is the \textbf{BEST} response?}
	{\item[\textbf{(A)}] Clarify that the business owner is accountable for the risk, with IT providing advice and support
		\item[\textbf{(B)}] Accept it because IT runs the servers
		\item[\textbf{(C)}] Refuse to engage
		\item[\textbf{(D)}] Assign it to the vendor}
	{A}
	{\item Risk ownership follows business process accountability.
		\item IT supports with expertise.
		\item Clear ownership improves decisions.}
	{Business processes carry business risk.}
	
	\hq{21}{Performance of Risk Owners}
	{Risk owners rarely update their risks. What is the \textbf{BEST} way to improve this?}
	{\item[\textbf{(A)}] Update the risks for them forever
		\item[\textbf{(B)}] Delete their risks
		\item[\textbf{(C)}] Define expectations, provide support and reminders, and include risk responsibilities in performance objectives
		\item[\textbf{(D)}] Hide the issue}
	{C}
	{\item Accountability needs expectations.
		\item Support makes it easier.
		\item Alignment with objectives gives it weight.}
	{Make risk ownership part of the job.}
	
	\hq{22}{Gaining Executive Support}
	{Executives view risk management as a compliance burden. What is the \textbf{BEST} approach?}
	{\item[\textbf{(A)}] Add more forms
		\item[\textbf{(B)}] Threaten with audits
		\item[\textbf{(C)}] Reduce reporting entirely
		\item[\textbf{(D)}] Demonstrate how risk information improves decisions and helps achieve strategic goals}
	{D}
	{\item Value in business terms earns support.
		\item Link risk to objectives.
		\item Share examples of decisions improved.}
	{Sell risk management as decision support.}
	
	\hq{23}{Open Communication of Risk}
	{A risk practitioner finds that important risks are not shared across departments. What should be improved?}
	{\item[\textbf{(A)}] Establish regular cross-functional risk communication and shared reporting channels
		\item[\textbf{(B)}] Keep all risks private
		\item[\textbf{(C)}] Share only good news
		\item[\textbf{(D)}] Limit sharing to IT}
	{A}
	{\item Interdependencies need shared awareness.
		\item Channels should be defined.
		\item Transparency builds trust.}
	{Risks do not respect department boundaries.}
	
	\hq{24}{Risk Culture Measurement}
	{How can the organization \textbf{BEST} assess its risk culture?}
	{\item[\textbf{(A)}] Through surveys, incident reporting patterns, and observed behavior of leaders and staff
		\item[\textbf{(B)}] By reading the policy manual
		\item[\textbf{(C)}] By counting servers
		\item[\textbf{(D)}] By asking only the board}
	{A}
	{\item Multiple indicators reveal culture.
		\item Reporting patterns show trust.
		\item Leaders' behavior is influential.}
	{Culture is observed in behavior.}
	
	\hq{25}{Strategy Pivot Risk}
	{Leadership announces a pivot to a new business model. What should the risk practitioner do?}
	{\item[\textbf{(A)}] Continue using last year's assessment unchanged
		\item[\textbf{(B)}] Oppose the pivot
		\item[\textbf{(C)}] Wait for incidents
		\item[\textbf{(D)}] Assess how the new strategy changes the risk profile and update appetite, scenarios, and reporting}
	{D}
	{\item Strategy changes risk exposure.
		\item Early assessment informs decisions.
		\item Reporting should reflect new risks.}
	{New strategy means a new risk profile.}
	
	\hq{26}{Independence of Risk Reporting}
	{A business head asks the risk function to soften a risk report before it goes to the board. What is the \textbf{BEST} response?}
	{\item[\textbf{(A)}] Soften the report as requested
		\item[\textbf{(B)}] Explain that the report must be accurate and offer to add the business's response and context
		\item[\textbf{(C)}] Remove the risk from the report
		\item[\textbf{(D)}] Send a different version to the business only}
	{B}
	{\item Boards need unbiased information.
		\item Business context can be added transparently.
		\item Altering facts damages credibility.}
	{Accurate reports, with context added openly.}
	
	% ============================================================
	\section{Risk Assessment (20 percent of the exam)}
	% ============================================================
	
	\hq{27}{Scenario Selection}
	{The team must choose risk scenarios for an assessment. Which approach is \textbf{BEST}?}
	{\item[\textbf{(A)}] Choose only scenarios that have already happened
		\item[\textbf{(B)}] Choose the most dramatic scenarios
		\item[\textbf{(C)}] Select plausible scenarios tied to key business processes, using input from owners and threat information
		\item[\textbf{(D)}] Choose scenarios at random}
	{C}
	{\item Scenarios should be relevant and plausible.
		\item Owners understand business impact.
		\item Threat data improves realism.}
	{Relevant scenarios beat dramatic ones.}
	
	\hq{28}{Estimating With Uncertainty}
	{Data on loss frequency is limited. What is the \textbf{BEST} way to estimate?}
	{\item[\textbf{(A)}] Use a single exact guess
		\item[\textbf{(B)}] Use calibrated ranges from experts and available data, and document assumptions
		\item[\textbf{(C)}] Refuse to estimate
		\item[\textbf{(D)}] Use the highest number possible}
	{B}
	{\item Ranges express uncertainty honestly.
		\item Calibration reduces bias.
		\item Assumptions can be reviewed and improved.}
	{Estimate with ranges and documented assumptions.}
	
	\hq{29}{Optimism Bias}
	{Assessors consistently rate their own area's risks as low. Which action \textbf{BEST} addresses this bias?}
	{\item[\textbf{(A)}] Apply independent challenge and calibration from the second line
		\item[\textbf{(B)}] Trust self-ratings completely
		\item[\textbf{(C)}] Eliminate self-assessments
		\item[\textbf{(D)}] Increase all ratings by one level}
	{A}
	{\item Independent challenge improves objectivity.
		\item Calibration aligns standards.
		\item Self-assessment is still useful with review.}
	{Challenge self-ratings independently.}
	
	\hq{30}{Interdependent Risks}
	{Two risks are rated medium separately, but both depend on the same data center. What should the assessment consider?}
	{\item[\textbf{(A)}] The common dependency, since a single failure could cause combined and higher impact
		\item[\textbf{(B)}] Nothing; rate each separately
		\item[\textbf{(C)}] Remove both risks
		\item[\textbf{(D)}] Lower both ratings}
	{A}
	{\item Dependencies create compounding effects.
		\item Aggregation reveals concentration.
		\item Shared causes call for shared responses.}
	{Look for common causes.}
	
	\hq{31}{Sensitivity Analysis}
	{What is the purpose of sensitivity analysis in risk assessment?}
	{\item[\textbf{(A)}] To hide uncertainty
		\item[\textbf{(B)}] To see how changes in assumptions affect results and identify the most influential inputs
		\item[\textbf{(C)}] To replace expert judgment
		\item[\textbf{(D)}] To increase the number of risks}
	{B}
	{\item It tests the robustness of estimates.
		\item Influential inputs deserve better data.
		\item It improves confidence in decisions.}
	{Test how much the answer depends on the assumptions.}
	
	\hq{32}{Monte Carlo Concept}
	{A quantitative model runs thousands of simulations using ranges for frequency and loss. What does this produce?}
	{\item[\textbf{(A)}] A distribution of possible outcomes showing the range and likelihood of losses
		\item[\textbf{(B)}] One guaranteed loss figure
		\item[\textbf{(C)}] A list of vulnerabilities
		\item[\textbf{(D)}] A compliance certificate}
	{A}
	{\item Simulation reflects uncertainty.
		\item Distributions support appetite comparisons.
		\item Quality depends on input quality.}
	{Simulation shows a range, not a promise.}
	
	\hq{33}{Data Quality for Assessment}
	{The assessment relies on an asset inventory known to be 40 percent out of date. What is the \textbf{BEST} action?}
	{\item[\textbf{(A)}] Disclose the limitation and improve inventory accuracy, prioritizing critical assets
		\item[\textbf{(B)}] Proceed and present results as precise
		\item[\textbf{(C)}] Cancel the assessment permanently
		\item[\textbf{(D)}] Use only the oldest data}
	{A}
	{\item Results depend on inputs.
		\item Transparency about limits preserves credibility.
		\item Prioritize critical assets first.}
	{Disclose data limits; do not hide them.}
	
	\hq{34}{Control Gap Analysis}
	{What is the \textbf{PRIMARY} output of a control gap analysis?}
	{\item[\textbf{(A)}] A list of employee names
		\item[\textbf{(B)}] Identification of differences between required or desired controls and those currently in place
		\item[\textbf{(C)}] A purchase order
		\item[\textbf{(D)}] A new policy only}
	{B}
	{\item Gaps show where risk may be unmitigated.
		\item They feed response planning.
		\item Required controls come from risk and obligations.}
	{A gap analysis shows what is missing.}
	
	\hq{35}{Red Team vs Pen Test}
	{How does a red team exercise differ from a standard penetration test?}
	{\item[\textbf{(A)}] It is the same thing
		\item[\textbf{(B)}] It simulates a realistic adversary with broader objectives and tests detection and response, not only vulnerabilities
		\item[\textbf{(C)}] It uses only automated scans
		\item[\textbf{(D)}] It avoids any realistic tactics}
	{B}
	{\item Red teams emulate real attackers.
		\item They test people, processes, and technology.
		\item Results inform detection improvements.}
	{Red teams test the defenders as well as the systems.}
	
	\hq{36}{Assessment Scope}
	{A pen test is requested. What should be agreed \textbf{FIRST}?}
	{\item[\textbf{(A)}] Scope, rules of engagement, and authorization, including safe handling of data
		\item[\textbf{(B)}] Which tools look best
		\item[\textbf{(C)}] The report's cover design
		\item[\textbf{(D)}] Nothing; just begin}
	{A}
	{\item Unauthorized testing creates legal risk.
		\item Rules protect production systems.
		\item Scope ensures relevance.}
	{Test with permission and boundaries.}
	
	\hq{37}{Legacy System Assessment}
	{A legacy system cannot be scanned safely. How should risk be assessed?}
	{\item[\textbf{(A)}] Treat the system as risk-free
		\item[\textbf{(B)}] Scan it anyway without planning
		\item[\textbf{(C)}] Use alternative methods such as architecture review, configuration review, and compensating control analysis
		\item[\textbf{(D)}] Ignore the system}
	{C}
	{\item Safe alternatives still provide insight.
		\item Compensating controls matter.
		\item Document limitations.}
	{If you cannot scan it, you still must assess it.}
	
	\hq{38}{Assessment Frequency}
	{How often should risk assessments of critical systems be performed?}
	{\item[\textbf{(A)}] At planned intervals based on risk and whenever significant changes occur
		\item[\textbf{(B)}] Once at system creation only
		\item[\textbf{(C)}] Only after incidents
		\item[\textbf{(D)}] Every ten years}
	{A}
	{\item Risk changes over time.
		\item Higher risk warrants more frequent review.
		\item Triggers ensure timeliness.}
	{Assess regularly and on change.}
	
	\hq{39}{Residual Risk Validation}
	{The assessed residual risk assumes a control works perfectly, but testing shows it works about half the time. What should happen?}
	{\item[\textbf{(A)}] Keep the original rating
		\item[\textbf{(B)}] Remove the control from the assessment
		\item[\textbf{(C)}] Re-rate residual risk based on observed control effectiveness
		\item[\textbf{(D)}] Cancel testing}
	{C}
	{\item Residual risk should reflect real performance.
		\item Testing evidence overrides assumptions.
		\item Update the register and inform owners.}
	{Residual risk follows actual control performance.}
	
	\hq{40}{Equal Scores}
	{Several risks have the same score. How should they be prioritized?}
	{\item[\textbf{(A)}] Pick alphabetically
		\item[\textbf{(B)}] Ignore all of them
		\item[\textbf{(C)}] Use the order received
		\item[\textbf{(D)}] Consider additional factors such as velocity, proximity, dependency, and strategic importance}
	{D}
	{\item Scores are only part of the picture.
		\item Timing and criticality matter.
		\item Judgment should be explained.}
	{Tie-breakers need business logic.}
	
	\hq{41}{New Vendor Assessment Depth}
	{How deep should a third-party risk assessment be?}
	{\item[\textbf{(A)}] Identical for every vendor
		\item[\textbf{(B)}] As minimal as possible
		\item[\textbf{(C)}] Based on vendor size only
		\item[\textbf{(D)}] Proportional to the vendor's criticality and the sensitivity of data and access involved}
	{D}
	{\item Risk-based depth saves effort.
		\item High-criticality vendors need more evidence.
		\item Reassess as relationships change.}
	{Depth follows criticality.}
	
	\hq{42}{IoT Assessment}
	{A factory adds many connected sensors. Which risk consideration is \textbf{MOST} important in the assessment?}
	{\item[\textbf{(A)}] Device inventory, firmware update ability, authentication, and network segmentation
		\item[\textbf{(B)}] Sensor color
		\item[\textbf{(C)}] Number of cables
		\item[\textbf{(D)}] Brand logos}
	{A}
	{\item IoT often has weak default security.
		\item Update capability determines long-term risk.
		\item Segmentation limits impact.}
	{IoT assessment starts with inventory and updateability.}
	
	\hq{43}{Expert Elicitation}
	{Experts disagree widely on a risk's likelihood. What is the \textbf{BEST} approach?}
	{\item[\textbf{(A)}] Use the most senior person's number
		\item[\textbf{(B)}] Discard the outliers always
		\item[\textbf{(C)}] Explore the reasoning behind each view and use structured methods to combine estimates
		\item[\textbf{(D)}] Use the lowest estimate}
	{C}
	{\item Disagreement reveals assumptions.
		\item Structured methods reduce bias.
		\item Document the rationale.}
	{Surface the reasoning behind the numbers.}
	
	\hq{44}{Business Dependency Mapping}
	{Why map dependencies between business processes and IT services?}
	{\item[\textbf{(A)}] To reduce printer use
		\item[\textbf{(B)}] To impress auditors
		\item[\textbf{(C)}] To understand how an IT failure would affect business processes and to prioritize protection
		\item[\textbf{(D)}] To avoid assessments}
	{C}
	{\item Mapping links IT risk to business impact.
		\item It informs BIA and recovery priorities.
		\item It reveals hidden single points of failure.}
	{Map the dependencies to understand real impact.}
	
	\hq{45}{Assessing Change}
	{A major system upgrade is planned. When should the risk assessment of the change be performed?}
	{\item[\textbf{(A)}] After go-live
		\item[\textbf{(B)}] During planning, before approval and implementation
		\item[\textbf{(C)}] One year later
		\item[\textbf{(D)}] Only if the upgrade fails}
	{B}
	{\item Early assessment informs go or no-go decisions.
		\item Mitigations can be built in.
		\item Late assessment is costlier.}
	{Assess change before approving it.}
	
	\hq{46}{Communicating Uncertainty}
	{How should the risk practitioner present estimates that carry considerable uncertainty?}
	{\item[\textbf{(A)}] Show a single precise number
		\item[\textbf{(B)}] Hide all uncertainty
		\item[\textbf{(C)}] Avoid presenting any estimate
		\item[\textbf{(D)}] Present ranges and confidence levels with clear assumptions, not false precision}
	{D}
	{\item Honest uncertainty builds trust.
		\item Decision makers can weigh ranges.
		\item False precision misleads.}
	{Do not claim more accuracy than you have.}
	
	% ============================================================
	\section{Risk Response and Reporting (32 percent of the exam)}
	% ============================================================
	
	\hq{47}{Response Selection Pressure}
	{Executives pressure the practitioner to choose the cheapest response for a high risk. What is the \textbf{BEST} approach?}
	{\item[\textbf{(A)}] Choose the cheapest regardless
		\item[\textbf{(B)}] Choose the costliest regardless
		\item[\textbf{(C)}] Avoid presenting options
		\item[\textbf{(D)}] Present response options with costs and residual risk so the owner can decide within appetite}
	{D}
	{\item Owners decide with full information.
		\item Residual risk shows what remains.
		\item Appetite is the benchmark.}
	{Options and residual risk enable decisions.}
	
	\hq{48}{Funding Action Plans}
	{Approved risk response plans lack funding. What should the practitioner do?}
	{\item[\textbf{(A)}] Escalate the funding gap and the resulting residual risk to decision makers
		\item[\textbf{(B)}] Close the plans
		\item[\textbf{(C)}] Reduce the risk rating
		\item[\textbf{(D)}] Wait silently}
	{A}
	{\item Unfunded plans leave risk unchanged.
		\item Decision makers must know the exposure.
		\item They may adjust funding or accept the risk.}
	{An unfunded plan is not a response.}
	
	\hq{49}{Designing KRIs}
	{Which criterion is \textbf{MOST} important when designing a KRI?}
	{\item[\textbf{(A)}] It is complicated to calculate
		\item[\textbf{(B)}] It is available only annually
		\item[\textbf{(C)}] It is linked to a specific risk and supports a decision when it changes
		\item[\textbf{(D)}] It reflects past events only}
	{C}
	{\item Relevance to a risk is essential.
		\item KRIs should trigger action.
		\item Timely and measurable data is needed.}
	{A KRI must lead to a decision.}
	
	\hq{50}{KCI Meaning}
	{A key control indicator shows that only 55 percent of required system backups succeeded last week. What does this \textbf{MOST} suggest?}
	{\item[\textbf{(A)}] Data loss risk has decreased
		\item[\textbf{(B)}] Backups are unnecessary
		\item[\textbf{(C)}] The indicator is irrelevant
		\item[\textbf{(D)}] The backup control is not operating effectively, increasing data loss risk}
	{D}
	{\item KCIs measure control performance.
		\item Low success means weaker protection.
		\item Investigate and remediate.}
	{A failing KCI signals a failing control.}
	
	\hq{51}{Dashboard Design}
	{An executive dashboard shows many technical metrics without context. How should it be improved?}
	{\item[\textbf{(A)}] Add more technical metrics
		\item[\textbf{(B)}] Remove all trends
		\item[\textbf{(C)}] Make it harder to read
		\item[\textbf{(D)}] Focus on key risks versus appetite, trends, and decisions needed}
	{D}
	{\item Executives need decision-ready information.
		\item Trends show direction.
		\item Context links metrics to appetite.}
	{Dashboards should drive decisions.}
	
	\hq{52}{Verifying Closure}
	{A team marks an audit finding as closed after implementing a fix. What should happen before the finding is truly closed?}
	{\item[\textbf{(A)}] Independent verification that the fix is effective
		\item[\textbf{(B)}] Nothing; the team reports closure
		\item[\textbf{(C)}] A new policy only
		\item[\textbf{(D)}] A vendor email}
	{A}
	{\item Self-reported closure may be incomplete.
		\item Verification provides assurance.
		\item Update residual risk afterward.}
	{Closed means verified.}
	
	\hq{53}{Exception Renewal}
	{A policy exception has been renewed five times without progress toward compliance. What should the practitioner do?}
	{\item[\textbf{(A)}] Renew it indefinitely
		\item[\textbf{(B)}] Quietly remove it
		\item[\textbf{(C)}] Escalate and require a realistic remediation plan or a formal risk acceptance at the right level
		\item[\textbf{(D)}] Hide the history}
	{C}
	{\item Repeated renewals signal unmanaged risk.
		\item Plans and dates drive progress.
		\item Acceptance needs proper authority.}
	{Exceptions should end, not renew forever.}
	
	\hq{54}{Sample Size}
	{A control operates thousands of times per month. How should testing be approached?}
	{\item[\textbf{(A)}] Test one instance only
		\item[\textbf{(B)}] Use a statistically or risk-based appropriate sample covering the period
		\item[\textbf{(C)}] Test only the first day
		\item[\textbf{(D)}] Test nothing}
	{B}
	{\item Samples must represent the population.
		\item Higher risk warrants larger samples.
		\item Document the approach.}
	{Sample enough to conclude.}
	
	\hq{55}{Continuous Controls Monitoring}
	{What is a key benefit of continuous controls monitoring?}
	{\item[\textbf{(A)}] It removes all risk
		\item[\textbf{(B)}] Faster detection of control failures through automated, ongoing testing
		\item[\textbf{(C)}] It eliminates control owners
		\item[\textbf{(D)}] It replaces management oversight}
	{B}
	{\item Automation tests frequently.
		\item Early detection shortens exposure.
		\item Owners still respond to alerts.}
	{Continuous tests catch failures sooner.}
	
	\hq{56}{Escalation Criteria}
	{The practitioner is unsure when to escalate risk issues. What should exist?}
	{\item[\textbf{(A)}] Personal discretion only
		\item[\textbf{(B)}] No escalation at all
		\item[\textbf{(C)}] Escalation only after incidents
		\item[\textbf{(D)}] Defined escalation criteria and thresholds approved by governance}
	{D}
	{\item Clear criteria ensure consistency.
		\item Thresholds link to appetite.
		\item Approval gives authority.}
	{Define when to escalate.}
	
	\hq{57}{Stale Register}
	{Half the risk register entries have not been reviewed in over a year. What is the \textbf{BEST} action?}
	{\item[\textbf{(A)}] Delete half the entries
		\item[\textbf{(B)}] Run a review cycle with owners, retire obsolete entries, and set recurring review dates
		\item[\textbf{(C)}] Ignore it
		\item[\textbf{(D)}] Rate all as low}
	{B}
	{\item Stale data misleads decisions.
		\item Owners confirm current status.
		\item Recurring dates maintain quality.}
	{A register is only as good as its freshness.}
	
	\hq{58}{Measuring Response Effectiveness}
	{How should the effectiveness of a risk response be measured?}
	{\item[\textbf{(A)}] By whether the project finished on time
		\item[\textbf{(B)}] By how much money was spent
		\item[\textbf{(C)}] By whether the actual risk exposure moved within appetite, as shown by indicators and testing
		\item[\textbf{(D)}] By the number of meetings}
	{C}
	{\item The goal is reduced exposure.
		\item Indicators and tests provide evidence.
		\item Activity is not the same as outcome.}
	{Effective response means reduced exposure.}
	
	\hq{59}{Residual Risk Sign-Off}
	{After mitigation, residual risk remains slightly above appetite. Who should decide whether to accept it?}
	{\item[\textbf{(A)}] The risk practitioner alone
		\item[\textbf{(B)}] The vendor
		\item[\textbf{(C)}] The internal auditor
		\item[\textbf{(D)}] The risk owner with senior management or governance approval as required}
	{D}
	{\item Owners and management hold authority.
		\item Above-appetite risk needs higher approval.
		\item Decisions must be documented.}
	{Above-appetite acceptance needs senior approval.}
	
	\hq{60}{Vendor Remediation Tracking}
	{A vendor promised to fix a significant weakness by the end of last quarter but has not. What should the practitioner do?}
	{\item[\textbf{(A)}] Wait without action
		\item[\textbf{(B)}] Remove the finding
		\item[\textbf{(C)}] Pay a bonus
		\item[\textbf{(D)}] Escalate under the contract, reassess the risk, and consider compensating controls or alternatives}
	{D}
	{\item Contracts provide escalation levers.
		\item Risk continues until fixed.
		\item Contingencies reduce exposure.}
	{Hold vendors to their remediation promises.}
	
	\hq{61}{Change and Risk}
	{Emergency changes bypass normal review frequently. What is the \textbf{BEST} control improvement?}
	{\item[\textbf{(A)}] Define an emergency change process with after-the-fact review and metrics on emergency use
		\item[\textbf{(B)}] Ban all emergency changes
		\item[\textbf{(C)}] Allow all changes without records
		\item[\textbf{(D)}] Ignore the pattern}
	{A}
	{\item Emergencies happen but need control.
		\item Post review ensures accountability.
		\item Metrics reveal abuse.}
	{Control the exception path too.}
	
	\hq{62}{DR Test Reporting}
	{A DR test met only half of its recovery time objectives. How should this be reported?}
	{\item[\textbf{(A)}] Report the gaps, the business impact, and a remediation plan with owners and dates
		\item[\textbf{(B)}] Report success anyway
		\item[\textbf{(C)}] Hide results
		\item[\textbf{(D)}] Blame the test team}
	{A}
	{\item Honest reporting enables improvement.
		\item Business impact gives context.
		\item Plans drive closure.}
	{Report failures as clearly as successes.}
	
	\hq{63}{Incident Metrics}
	{Which incident metric is \textbf{MOST} useful for management oversight?}
	{\item[\textbf{(A)}] Number of emails about incidents
		\item[\textbf{(B)}] Trends in severity, time to detect, and time to contain or recover
		\item[\textbf{(C)}] Names of those involved only
		\item[\textbf{(D)}] The color of the incident form}
	{B}
	{\item Trends show whether capability is improving.
		\item Time metrics show responsiveness.
		\item Severity reflects business impact.}
	{Trend the outcomes.}
	
	\hq{64}{Insurance Exclusions}
	{Before relying on cyber insurance as a response, what should be reviewed?}
	{\item[\textbf{(A)}] The insurer's logo
		\item[\textbf{(B)}] Only the premium
		\item[\textbf{(C)}] Nothing
		\item[\textbf{(D)}] Coverage limits, exclusions, and conditions, such as required controls}
	{D}
	{\item Policies may exclude events or require controls.
		\item Limits may be lower than potential loss.
		\item Misunderstanding coverage creates false assurance.}
	{Read the exclusions before you rely on the policy.}
	
	\hq{65}{Validating Compensating Controls}
	{A compensating control is used in place of a standard control. What should be done?}
	{\item[\textbf{(A)}] Assume it is equivalent
		\item[\textbf{(B)}] Assess that it addresses the same risk to an equivalent degree and document it
		\item[\textbf{(C)}] Remove the standard
		\item[\textbf{(D)}] Keep it secret}
	{B}
	{\item Compensating controls must be effective.
		\item Documentation supports review.
		\item Test periodically.}
	{Compensate only with equivalence.}
	
	\hq{66}{System Decommissioning}
	{A system is being retired. Which risk activity is \textbf{MOST} important?}
	{\item[\textbf{(A)}] Turn it off and forget it
		\item[\textbf{(B)}] Leave accounts active
		\item[\textbf{(C)}] Keep it online unmonitored
		\item[\textbf{(D)}] Ensure data is migrated or securely destroyed, access removed, and records updated}
	{D}
	{\item Retired systems can leak data.
		\item Orphan accounts are risks.
		\item Inventory must reflect reality.}
	{Decommission thoroughly.}
	
	\hq{67}{Patch Exceptions}
	{A critical patch cannot be applied for three months due to business constraints. What is the \textbf{BEST} approach?}
	{\item[\textbf{(A)}] Ignore the patch
		\item[\textbf{(B)}] Remove monitoring
		\item[\textbf{(C)}] Delay reporting
		\item[\textbf{(D)}] Record a time-bound exception with compensating controls and owner approval}
	{D}
	{\item Exceptions need compensating controls.
		\item Time limits keep pressure on closure.
		\item Owner approval maintains accountability.}
	{Delay with controls, not silence.}
	
	\hq{68}{Risk Trending}
	{Why is trend analysis useful in risk reporting?}
	{\item[\textbf{(A)}] It hides outliers
		\item[\textbf{(B)}] It shows whether exposure is improving or worsening and supports proactive decisions
		\item[\textbf{(C)}] It replaces appetite
		\item[\textbf{(D)}] It is only for auditors}
	{B}
	{\item Direction matters as much as level.
		\item Trends can predict threshold breaches.
		\item They support timely action.}
	{Report the direction as well as the position.}
	
	\hq{69}{Regulator Reporting Accuracy}
	{Data for a regulatory report is compiled manually from several sources. What is the \textbf{BEST} improvement?}
	{\item[\textbf{(A)}] Submit unreviewed
		\item[\textbf{(B)}] Rely on one person's memory
		\item[\textbf{(C)}] Define controls over data accuracy, reconciliation, and review before submission
		\item[\textbf{(D)}] Skip reconciliation}
	{C}
	{\item Regulatory errors carry penalties.
		\item Reconciliation verifies completeness.
		\item Review adds assurance.}
	{Control the data behind regulatory reports.}
	
	\hq{70}{Communicating to Non-Technical Owners}
	{How should a risk practitioner explain a technical vulnerability to a business owner?}
	{\item[\textbf{(A)}] Translate it into potential business impact, likelihood, and available response options
		\item[\textbf{(B)}] Use detailed exploit code
		\item[\textbf{(C)}] Use only technical jargon
		\item[\textbf{(D)}] Avoid explaining}
	{A}
	{\item Owners decide in business terms.
		\item Options enable choice.
		\item Clear language builds trust.}
	{Translate technical risk into business impact.}
	
	\hq{71}{Control Optimization}
	{Several overlapping controls address the same low risk at high cost. What should be done?}
	{\item[\textbf{(A)}] Evaluate rationalizing redundant controls while keeping risk within appetite
		\item[\textbf{(B)}] Add more controls
		\item[\textbf{(C)}] Remove all controls
		\item[\textbf{(D)}] Ignore costs}
	{A}
	{\item Efficiency matters in risk programs.
		\item Redundancy is valuable only where justified.
		\item Appetite is the constraint.}
	{Optimize controls, do not just accumulate them.}
	
	\hq{72}{Reliance on Others' Assurance}
	{May the risk practitioner rely on an internal audit report as evidence about a control?}
	{\item[\textbf{(A)}] Never use others' work
		\item[\textbf{(B)}] Yes, after considering the scope, timing, independence, and quality of the audit work
		\item[\textbf{(C)}] Use any report blindly
		\item[\textbf{(D)}] Use only vendor claims}
	{B}
	{\item Reliance requires judgment.
		\item Scope must match the control.
		\item Timing affects relevance.}
	{Rely on others' work only after evaluating it.}
	
	\hq{73}{KRI Data Reliability}
	{A KRI relies on data that is manually entered and inconsistent. What is the \textbf{BEST} action?}
	{\item[\textbf{(A)}] Ignore the KRI
		\item[\textbf{(B)}] Improve data collection, automate where possible, and validate quality
		\item[\textbf{(C)}] Use it as is
		\item[\textbf{(D)}] Hide the data source}
	{B}
	{\item Unreliable data misleads.
		\item Automation increases consistency.
		\item Validation builds confidence.}
	{Fix the data before trusting the indicator.}
	
	\hq{74}{Risk Acceptance Review}
	{A risk was accepted two years ago. When should the acceptance be reviewed?}
	{\item[\textbf{(A)}] Periodically and whenever conditions change significantly
		\item[\textbf{(B)}] Never
		\item[\textbf{(C)}] Only if someone complains
		\item[\textbf{(D)}] After the owner leaves only}
	{A}
	{\item Acceptance is based on conditions that may change.
		\item Reviews confirm continued validity.
		\item Owner changes also trigger reviews.}
	{Acceptance is not forever.}
	
	\hq{75}{Third-Party Reporting}
	{How should third-party risk be reported to management?}
	{\item[\textbf{(A)}] List every vendor equally
		\item[\textbf{(B)}] Highlight critical vendors, key findings, concentration, and remediation status against appetite
		\item[\textbf{(C)}] Report only vendor names
		\item[\textbf{(D)}] Skip reporting}
	{B}
	{\item Focus on what matters.
		\item Concentration is a key concern.
		\item Status shows progress.}
	{Report the vendors that matter.}
	
	\hq{76}{Final Risk Position}
	{Management asks for the overall risk posture after a series of responses. What should the practitioner provide?}
	{\item[\textbf{(A)}] Only the list of completed tasks
		\item[\textbf{(B)}] The number of hours worked
		\item[\textbf{(C)}] Updated residual risk positions against appetite, with trends and open decisions
		\item[\textbf{(D)}] A guarantee of no incidents}
	{C}
	{\item Posture means exposure relative to appetite.
		\item Trends and open items guide decisions.
		\item No one can guarantee zero incidents.}
	{Report posture, not just progress.}
	
	\hq{77}{Closing the Loop With Owners}
	{After a risk report is issued, the risk owners do not respond. What should the practitioner do?}
	{\item[\textbf{(A)}] Assume silence means acceptance
		\item[\textbf{(B)}] Close the risks
		\item[\textbf{(C)}] Stop reporting
		\item[\textbf{(D)}] Follow up, confirm they understood the risks and options, and escalate if no decision is made}
	{D}
	{\item Silence is not acceptance.
		\item Follow-up ensures understanding.
		\item Escalation secures a decision.}
	{Silence is not a decision.}
	
	\hq{78}{Risk Appetite Breach Response}
	{A KRI shows that a risk has breached its tolerance threshold for the second month. What is the \textbf{BEST} action?}
	{\item[\textbf{(A)}] Escalate to the risk owner and governance per the defined process and agree on corrective actions
		\item[\textbf{(B)}] Wait one more month
		\item[\textbf{(C)}] Change the threshold to remove the breach
		\item[\textbf{(D)}] Delete the KRI}
	{A}
	{\item Persistent breaches need action.
		\item Thresholds reflect approved tolerance.
		\item Governance decides on exceptions.}
	{Breaches trigger escalation, not threshold changes.}
	
	% ============================================================
	\section{Information Technology and Security (22 percent of the exam)}
	% ============================================================
	
	\hq{79}{Email Authentication}
	{Attackers keep sending emails that spoof the company's domain. Which control set \textbf{BEST} reduces this?}
	{\item[\textbf{(A)}] SPF, DKIM, and DMARC email authentication with enforcement
		\item[\textbf{(B)}] Longer email subject lines
		\item[\textbf{(C)}] Larger mailboxes
		\item[\textbf{(D)}] More email signatures}
	{A}
	{\item These standards verify sender legitimacy.
		\item Enforcement tells receivers to reject spoofed mail.
		\item Monitoring reports reveal abuse.}
	{Authenticate your domain's email.}
	
	\hq{80}{Endpoint Detection}
	{Traditional antivirus misses several advanced attacks. Which capability \textbf{BEST} improves detection on endpoints?}
	{\item[\textbf{(A)}] A bigger monitor
		\item[\textbf{(B)}] Password sticky notes
		\item[\textbf{(C)}] Endpoint detection and response with behavioral monitoring
		\item[\textbf{(D)}] Disabling updates}
	{C}
	{\item Behavior-based tools detect unknown threats.
		\item EDR supports investigation and containment.
		\item It complements, not replaces, other controls.}
	{Detect behavior, not just signatures.}
	
	\hq{81}{Data Loss Prevention}
	{What is the main purpose of data loss prevention tools?}
	{\item[\textbf{(A)}] To speed up networks
		\item[\textbf{(B)}] To compress files
		\item[\textbf{(C)}] To detect and block unauthorized transfer of sensitive data based on policy
		\item[\textbf{(D)}] To replace backups}
	{C}
	{\item DLP enforces data handling rules.
		\item Classification improves accuracy.
		\item Tuning reduces false positives.}
	{DLP guards the exits.}
	
	\hq{82}{Encryption at Rest}
	{A laptop containing sensitive data is lost. Which control \textbf{BEST} limits the impact?}
	{\item[\textbf{(A)}] A screen saver
		\item[\textbf{(B)}] A louder alarm
		\item[\textbf{(C)}] A longer warranty
		\item[\textbf{(D)}] Full-disk encryption with strong key protection}
	{D}
	{\item Encrypted data is unreadable without keys.
		\item Key protection is critical.
		\item Remote wipe adds further protection.}
	{Lost device, protected data.}
	
	\hq{83}{Transport Security}
	{Sensitive data is sent between systems over an untrusted network. Which control \textbf{BEST} protects it in transit?}
	{\item[\textbf{(A)}] Plain text transfer
		\item[\textbf{(B)}] Obscure port numbers
		\item[\textbf{(C)}] Compression only
		\item[\textbf{(D)}] Strong transport encryption such as current TLS versions}
	{D}
	{\item Encryption prevents eavesdropping.
		\item Current protocol versions avoid known weaknesses.
		\item Certificate validation prevents impersonation.}
	{Encrypt in transit with current protocols.}
	
	\hq{84}{Single Sign-On}
	{What is a security benefit of single sign-on with strong authentication?}
	{\item[\textbf{(A)}] No need for any authentication
		\item[\textbf{(B)}] Elimination of all risk
		\item[\textbf{(C)}] Fewer passwords to manage and centralized control of access and logging
		\item[\textbf{(D)}] Shared passwords among staff}
	{C}
	{\item Central policies are easier to enforce.
		\item Users avoid password reuse.
		\item The identity provider becomes critical and must be protected.}
	{SSO simplifies control but concentrates risk.}
	
	\hq{85}{Role-Based Access}
	{Access permissions are assigned individually and have become inconsistent. What is the \textbf{BEST} improvement?}
	{\item[\textbf{(A)}] Give everyone admin rights
		\item[\textbf{(B)}] Let managers use any approach
		\item[\textbf{(C)}] Remove all reviews
		\item[\textbf{(D)}] Implement role-based access control aligned to job functions and review it periodically}
	{D}
	{\item Roles standardize permissions.
		\item Reviews catch drift.
		\item Least privilege is easier to maintain.}
	{Roles make least privilege manageable.}
	
	\hq{86}{Database Protection}
	{Which combination \textbf{BEST} protects a database containing sensitive records?}
	{\item[\textbf{(A)}] A single shared account
		\item[\textbf{(B)}] Public internet exposure
		\item[\textbf{(C)}] Access controls, encryption, activity monitoring, and secure configuration
		\item[\textbf{(D)}] No logging}
	{C}
	{\item Layered controls reduce exposure.
		\item Monitoring detects misuse.
		\item Configuration prevents known weaknesses.}
	{Protect data where it lives.}
	
	\hq{87}{Backup Strategy}
	{Which approach is the \textbf{BEST} backup practice for resilience against ransomware?}
	{\item[\textbf{(A)}] Multiple copies on different media with at least one offline or immutable copy, regularly tested
		\item[\textbf{(B)}] One copy on the same server
		\item[\textbf{(C)}] Backups never tested
		\item[\textbf{(D)}] Backups on a mapped network drive only}
	{A}
	{\item Diversity reduces common failure.
		\item Offline copies resist encryption attacks.
		\item Testing proves recoverability.}
	{Multiple copies, one offline, all tested.}
	
	\hq{88}{Web Application Firewall}
	{What does a web application firewall primarily help protect against?}
	{\item[\textbf{(A)}] Disk failures
		\item[\textbf{(B)}] Printer jams
		\item[\textbf{(C)}] Power outages
		\item[\textbf{(D)}] Common web attacks such as injection and cross-site scripting targeting web applications}
	{D}
	{\item A WAF filters malicious application traffic.
		\item It complements secure coding.
		\item Rules need tuning.}
	{WAF protects web apps, not hardware.}
	
	\hq{89}{Authenticated Scanning}
	{Why is authenticated vulnerability scanning often preferred?}
	{\item[\textbf{(A)}] It is always faster
		\item[\textbf{(B)}] It avoids all false positives
		\item[\textbf{(C)}] It requires no credentials
		\item[\textbf{(D)}] It provides deeper, more accurate results by inspecting configurations and installed software}
	{D}
	{\item Credentials allow internal visibility.
		\item Fewer false negatives result.
		\item Credentials must be protected.}
	{Authenticated scans see more.}
	
	\hq{90}{Time Synchronization}
	{Logs from different systems cannot be correlated because their clocks differ. What is the \textbf{BEST} fix?}
	{\item[\textbf{(A)}] Ignore times in logs
		\item[\textbf{(B)}] Turn off logging
		\item[\textbf{(C)}] Synchronize systems to a reliable time source
		\item[\textbf{(D)}] Use different time zones deliberately}
	{C}
	{\item Accurate timestamps are essential for investigation.
		\item Correlation depends on consistent time.
		\item Time services must be protected.}
	{Correlation requires consistent clocks.}
	
	\hq{91}{Cloud Identity Risk}
	{Cloud administrator accounts use only passwords. What is the \textbf{BEST} recommendation?}
	{\item[\textbf{(A)}] Share the admin password
		\item[\textbf{(B)}] Require multi-factor authentication, least privilege, and monitoring for administrative access
		\item[\textbf{(C)}] Use long-lived keys in code
		\item[\textbf{(D)}] Remove logging}
	{B}
	{\item Admin accounts are high-value targets.
		\item MFA reduces compromise.
		\item Monitoring detects misuse.}
	{Protect cloud admins most.}
	
	\hq{92}{Secrets in Code}
	{Developers store passwords and API keys in source code. What is the \textbf{BEST} remedy?}
	{\item[\textbf{(A)}] Keep secrets in code but hide the repository name
		\item[\textbf{(B)}] Share secrets by chat
		\item[\textbf{(C)}] Never rotate keys
		\item[\textbf{(D)}] Use a secrets management solution, rotate exposed secrets, and scan code for credentials}
	{D}
	{\item Secrets in code leak easily.
		\item Rotation limits exposure.
		\item Scanning prevents recurrence.}
	{Keep secrets out of code.}
	
	\hq{93}{Containers}
	{A company runs applications in containers. Which practice is \textbf{MOST} important for managing risk?}
	{\item[\textbf{(A)}] Pull images from any source
		\item[\textbf{(B)}] Use trusted images, scan for vulnerabilities, limit privileges, and secure orchestration
		\item[\textbf{(C)}] Run all containers as root
		\item[\textbf{(D)}] Ignore updates}
	{B}
	{\item Image provenance reduces supply chain risk.
		\item Least privilege contains compromise.
		\item Orchestration security is crucial.}
	{Trust, scan, limit.}
	
	\hq{94}{Serverless Risk}
	{Which responsibility remains with the customer when using serverless functions?}
	{\item[\textbf{(A)}] Secure code, permissions, and handling of data
		\item[\textbf{(B)}] Physical servers
		\item[\textbf{(C)}] Hypervisor patching
		\item[\textbf{(D)}] Data center cooling}
	{A}
	{\item Providers manage infrastructure.
		\item Customers control code and access.
		\item Misconfigured permissions are common issues.}
	{Serverless still has code and permission risk.}
	
	\hq{95}{Privacy by Design}
	{What does privacy by design mean?}
	{\item[\textbf{(A)}] Building privacy protections into systems and processes from the start
		\item[\textbf{(B)}] Adding privacy after launch
		\item[\textbf{(C)}] Collecting maximum data
		\item[\textbf{(D)}] Avoiding documentation}
	{A}
	{\item Early design is more effective.
		\item It supports minimization and purpose limits.
		\item It reduces compliance costs.}
	{Design privacy in, not on.}
	
	\hq{96}{Pseudonymization}
	{How does pseudonymization differ from anonymization?}
	{\item[\textbf{(A)}] They are identical
		\item[\textbf{(B)}] Pseudonymized data can be re-identified with additional information, whereas properly anonymized data cannot reasonably be re-identified
		\item[\textbf{(C)}] Anonymization is reversible by anyone
		\item[\textbf{(D)}] Pseudonymization removes all data}
	{B}
	{\item Pseudonyms reduce but do not remove privacy risk.
		\item Additional information must be protected.
		\item True anonymization is difficult.}
	{Pseudonymized is not anonymous.}
	
	\hq{97}{Chain of Custody}
	{Why is chain of custody important in digital forensics?}
	{\item[\textbf{(A)}] It speeds up analysis
		\item[\textbf{(B)}] It lowers costs
		\item[\textbf{(C)}] It replaces backups
		\item[\textbf{(D)}] It documents who handled evidence and when, preserving its integrity and admissibility}
	{D}
	{\item Documentation shows evidence was not altered.
		\item Legal proceedings may depend on it.
		\item Procedures should be followed strictly.}
	{Evidence needs a documented trail.}
	
	\hq{98}{Capacity Management}
	{Which risk does poor capacity management \textbf{MOST} directly create?}
	{\item[\textbf{(A)}] Excess privacy
		\item[\textbf{(B)}] Performance degradation or outages when demand exceeds resources
		\item[\textbf{(C)}] Better compliance
		\item[\textbf{(D)}] Faster recovery}
	{B}
	{\item Availability depends on adequate capacity.
		\item Trending usage supports planning.
		\item Cloud elasticity helps but needs oversight.}
	{Plan capacity before demand outruns it.}
	
	\hq{99}{IT Asset Disposal}
	{Retired laptops and drives are being sent for recycling. What is the \textbf{BEST} control?}
	{\item[\textbf{(A)}] Delete visible files only
		\item[\textbf{(B)}] Securely wipe or destroy storage media and keep disposal records
		\item[\textbf{(C)}] Donate unwiped
		\item[\textbf{(D)}] Store in an unlocked room}
	{B}
	{\item Simple deletion leaves data recoverable.
		\item Records provide proof.
		\item Vendors should be verified.}
	{Sanitize before disposal.}
	
	\hq{100}{Security Misconfiguration Management}
	{Cloud and server configurations frequently drift from approved settings. What is the \textbf{BEST} control?}
	{\item[\textbf{(A)}] Review configurations once every five years
		\item[\textbf{(B)}] Trust administrators' memory
		\item[\textbf{(C)}] Define secure baselines and use automated configuration monitoring with remediation workflows
		\item[\textbf{(D)}] Disable monitoring}
	{C}
	{\item Drift is a common cause of exposure.
		\item Automation detects deviations quickly.
		\item Workflows ensure fixes.}
	{Baseline it, monitor it, fix the drift.}
	
	% ============================================================
	\section{Conclusion}
	% ============================================================
	
	\begin{frame}{Exam Mindset Summary}
		\begin{block}{Remember These Core Principles}
			\begin{itemize}
				\item \textbf{Business First:} Risk is defined by its effect on business objectives.
				\item \textbf{Owners Decide:} Practitioners inform; owners accept, mitigate, transfer, or avoid.
				\item \textbf{Three Lines:} The first line owns risk, the second oversees, and the third gives independent assurance.
				\item \textbf{Ethics Matter:} Honesty, independence, and confidentiality underpin credible risk advice.
				\item \textbf{Evidence Over Opinion:} Use data, scenarios, and calibrated ratings.
				\item \textbf{Proportionality:} Controls should fit the risk and its cost.
				\item \textbf{Accountability Cannot be Outsourced:} Vendors and cloud providers share duties, but accountability remains with the organization.
			\end{itemize}
		\end{block}
	\end{frame}
	
	\begin{frame}{Thank You}
		\begin{center}
			{\Huge \textbf{Good Luck on Your CRISC Exam!}}
			
			\vspace{1cm}
			
			\normalsize
			Real exam questions test your ability to apply these concepts in practical, real-world scenarios.
		\end{center}
	\end{frame}
	
\end{document}